\chapter{Industrial Context and Internship Mission}
\label{chap:context}

\section{OCP Jorf Lasfar and the project setting}

OCP Group operates an integrated phosphate value chain, from extraction to transformation and plant-nutrition products. At Jorf Lasfar, production units interact with acid and fertilizer facilities, utilities, storage and port logistics \cite{ocpwho,ocpindustrial}. For this internship, the relevant point is not the corporate scale itself. It is the need to coordinate material, energy, equipment condition and quality information across a continuous industrial system.

\begin{figure}[htbp]
\centering
\begin{subfigure}[t]{.48\textwidth}
  \centering
  \includegraphics[width=\linewidth,height=4.8cm,keepaspectratio]{jorf_lasfar_industrial_platform.png}
  \caption{Integrated processing platform.}
\end{subfigure}\hfill
\begin{subfigure}[t]{.48\textwidth}
  \centering
  \includegraphics[width=\linewidth,height=4.8cm,keepaspectratio]{jorf_lasfar_port.png}
  \caption{Port and logistics interface.}
\end{subfigure}
\caption{Jorf Lasfar industrial context.}
\source{Supplied ENSAM/OCP organization report \cite{orgreport}; contextual images, not evidence of prototype deployment.}
\label{fig:jorf-context}
\end{figure}

My fieldwork focused on the soluble MAP installation. I followed the production sequence from reactant preparation and neutralization to crystallization, centrifugal separation, drying, cooling and handling. The PCS 7 screens provided a useful view of how these stages are supervised. They informed the process analysis in this report, but the prototype has no live PCS 7 connection and does not reproduce its control logic.

\section{Observed supervision problem}

Process variables are available continuously in the control environment, whereas final moisture, density and product temperature are confirmed by periodic laboratory analysis. The last laboratory value is authoritative for the sampled material, but it becomes progressively older while the process continues to change. This is the gap addressed by the project.

The useful digital response is not to hide that age. The interface must state whether a number is a current process observation, a laboratory measurement, a model estimate or a diagnostic interpretation. Timestamp, data freshness and model version are therefore part of the functional design, not optional technical metadata.

\section{Mission, contribution and boundary}

The mission was to design and demonstrate an intelligent supervision architecture around dryer moisture and abnormal-operation visibility. I carried out the process-to-data translation, causal time alignment, feature design, model comparison, anomaly interpretation, database integration, runtime checks and dashboard information design. Python performs the analytical work, PostgreSQL provides traceable persistence, and Power BI presents the result.

During the internship, I built the prototype layer around the existing process, PCS 7 supervision context and laboratory practice. My work covered the process-to-data translation, model workflow, traceable storage and operator-facing dashboard. A convincing interface should still not be mistaken for a deployed plant system.

The deliverable remains a controlled proof of concept. It does not write to equipment, automate conformity decisions, replace laboratory control or claim industrial-scale accuracy. Future use would require approved historian access, tag and clock reconciliation, representative campaigns, cybersecurity review, shadow validation and operating procedures jointly owned by OCP stakeholders.

\begin{engineeringbox}[Design requirements taken from the mission]
Every displayed value retains a timestamp and provenance; stale or incomplete states remain visible; anomaly evidence leads to a verification prompt; model and feature versions are traceable; and the complete analytical layer remains outside the control loop.
\end{engineeringbox}

The following chapter develops the process basis behind these requirements, with particular attention to the dryer, its material delay and the measurements available to the prototype.
